\documentclass[12pt,a4paper]{article}

\input{/home/tabaka/Documents/GitHub/Defs/defs.tex}

\setcounter{zadaniaRozdzial}{3}

\begin{document}
	
	\begin{center}
		\LARGE Pola czworokątów - równoległoboki
	\end{center}
	
	\ZbiorZadanie{Boki równoległoboku mają długości $6\text{ cm}$ i $10\text{ cm}$, a jego pole jest równe $90\text{ cm}^2$. Oblicz różnicę długości wysokości tego równoległoboku.}
	
	\ZbiorZadanie{Oblicz pole równoległoboku, którego boki mają długości $8\text{ cm}$ i $12\text{ cm}$, a kąt:
		
		\begin{enumerate}[a)]
			\item ostry ma miarę $45^\circ$,
			\item rozwarty jest pięć razy większy od kąta ostrego.
	\end{enumerate}}
	
	\ZbiorZadanie{Równoległobok o bokach długości $6\text{ cm}$ i $12\text{ cm}$ ma pole $P$. Oblicz miarę kąta ostrego tego równoległoboku, gdy:
		
		\begin{enumerate}[a)]
			\item $P=36\text{ cm}^2$,
			\item $P=36\sqrt{2}\text{ cm}^2$,
			\item $P=36\sqrt{3}\text{ cm}^2$.
	\end{enumerate}}
	
	\ZbiorZadanie{Uzasadnij, że przekątne równoległoboku dzielą go na cztery trójkąty o równych polach.}
	
	\ZbiorZadanie{Oblicz pole równoległoboku, którego przekątne mają długości $5\text{ cm}$ i $10\text{ cm}$ oraz przecinają się pod kątem $30^\circ$.}
	
	\ZbiorZadanie{Oblicz pole i wysokość rombu:
		
		\begin{enumerate}[a)]
			\item o boku długości $13\text{ cm}$ i jednej przekątnej równej $24\text{ cm}$,
			\item o boku długości $6\text{ cm}$ i kącie ostrym $\alpha$ takim, że $\cos\alpha=\frac{1}{3}$.
	\end{enumerate}}
	
	\ZbiorZadanie{Oblicz pole równoległoboku, którego kąt rozwarty ma miarę $150^\circ$, a boki mają długości 4 cm i 9 cm.}
	
	\ZbiorZadanie{Pole równoległoboku wynosi $8\text{ cm}^2$, obwód $24\text{ cm}$, a sinus kąta ostrego jest równy $0{,}25$. Oblicz wysokości tego równoległoboku.}
	
	\ZbiorZadanie{Przekątne równoległoboku o długościach $6\text{ cm}$ i $10\text{ cm}$ tworzą z jednym z boków odpowiednio kąty $12^\circ$ i $33^\circ$. Oblicz pole tego równoległoboku.}
	
	\ZbiorZadanie{Oblicz sinus kąta ostrego rombu o przekątnych długości 12 cm i 16 cm.}
	
	\ZbiorZadanie{Oblicz wysokość rombu o polu równym $120\text{ cm}^2$ i obwodzie  cm.}
	
	\ZbiorZadanie{W rombie o boku długości $4\text{ cm}$ i kącie ostrym $60^\circ$ połączono odcinkami środki sąsiednich boków i otrzymano prostokąt. Oblicz jego pole.}
	
	\newpage
	
	\begin{center}
		\LARGE Pola czworokątów - trapezy
	\end{center}
	
	\ZbiorZadanie{Oblicz pole i obwód trapezu równoramiennego, którego:
		
		\begin{enumerate}[a)]
			\item podstawy mają długości $10\text{ cm}$ i $6\text{ cm}$, a przekątna ma długość $9\text{ cm}$,
			\item ramię ma długość $7\text{ cm}$, krótsza podstawa ma długość $5\text{ cm}$, a wysokość jest równa $3\sqrt{5}\text{ cm}$,
			\item krótsza podstawa ma długość $3\sqrt{3}\text{ cm}$, dłuższa podstawa ma długość $7\sqrt{3}\text{ cm}$, a miara kąta rozwartego wynosi $150^\circ$,
			\item jedna z podstaw jest trzy razy dłuższa od drugiej, przekątna ma długość $9\text{ cm}$, a wysokość jest równa $3\sqrt{5}\text{ cm}$.
	\end{enumerate}}
	
	\ZbiorZadanie{Dany jest trapez o wysokości $4\text{ cm}$, krótszej podstawie $5\sqrt{5}\text{ cm}$ oraz kątach ostrych $\alpha$ i $\beta$ takich, że $\sin\alpha=\frac{2}{3}$ i $\sin\beta=\frac{4\sqrt{21}}{21}$. Oblicz pole tego trapezu.}
	
	\ZbiorZadanie{Boki deltoidu mają długości $3\sqrt{2}\text{ cm}$ i $5\text{ cm}$, a krótsza przekątna ma długość $6\text{ cm}$. Oblicz pole tego deltoidu.}
	
	\ZbiorZadanie{W trapezie prostokątnym o polu $90\text{ cm}^2$ i kącie ostrym $45^\circ$ dłuższa przekątna tworzy z podstawami kąt $\alpha$ taki, że $\tg\alpha=\frac{1}{3}$. Oblicz obwód tego trapezu.}
	
	\ZbiorZadanie{Trapez równoramienny o podstawach długości $2\text{ cm}$ i $4\text{ cm}$ ma pole równe $18\text{ cm}^2$. Oblicz sinus kąta, pod jakim przecinają się przekątne tego trapezu.}
	
	\ZbiorZadanie{Działkę budowlaną w kształcie trapezu o kolejnych bokach długości $50\text{ m}$, $25\text{ m}$, $20\text{ m}$ i $25\text{ m}$ podzielono linią równoległą do podstaw tak, że obwody nowo powstałych działek są równe. Ile metrów bieżących siatki potrzeba do ogrodzenia obu działek (siatka między działkami jest wspólna), jeżeli będą one miały furtki o szerokości $1{,}5\text{ m}$? Oblicz pola tych działek.}
	
	\ZbiorZadanie{Przekątne $AC$ i $BD$ trapezu $ABCD$ o podstawach $AB$ i $CD$ przecinają się w punkcie $E$. Wykaż, że trójkąty $AED$ i $BEC$ mają równe pola.}
	
	\ZbiorZadanie{Udowodnij, że odcinek łączący środki ramion $AD$ i $BC$ trapezu $ABCD$ ma długość $\frac{a+b}{2}$.}
	
	\ZbiorZadanie{Ramiona trapezu mają długości $3\text{ cm}$ i $4\text{ cm}$, krótsza podstawa ma $7{,}5\text{ cm}$, a odcinek łączący środki ramion ma długość $10\text{ cm}$. Oblicz długość dłuższej podstawy tego trapezu i jego pole.}
	
	\ZbiorZadanie{Boki deltoidu mają długości $5\text{ cm}$ i $12\text{ cm}$, a dłuższa przekątna ma długość $13\text{ cm}$. Oblicz pole tego deltoidu.}
	
	\ZbiorZadanie{Oblicz obwód i pole trapezu równoramiennego o krótszej podstawie 12 cm, kącie ostrym $30^\circ$ i wysokości 4 cm.}

	
\end{document}